\documentclass[11pt]{article}
\usepackage[T1]{fontenc}
\usepackage{lmodern}
\usepackage[letterpaper,margin=1in]{geometry}
\usepackage{microtype}
\usepackage{amsmath,amssymb,amsthm,mathtools,mathrsfs}
\usepackage{enumitem}
\usepackage[hidelinks]{hyperref}
\numberwithin{equation}{section}
\newtheorem{theorem}{Theorem}[section]
\newtheorem{proposition}[theorem]{Proposition}
\newtheorem{lemma}[theorem]{Lemma}
\newtheorem{corollary}[theorem]{Corollary}
\theoremstyle{definition}
\newtheorem{definition}[theorem]{Definition}
\theoremstyle{remark}
\newtheorem{remark}[theorem]{Remark}
\newcommand{\C}{\mathbb C}
\newcommand{\R}{\mathbb R}
\newcommand{\Z}{\mathbb Z}
\newcommand{\Gm}{\mathbb G_m}
\newcommand{\hX}{\mathfrak h_X}
\newcommand{\hY}{\mathfrak h_Y}
\newcommand{\dd}{\,\mathrm d}
\DeclareMathOperator{\SL}{SL}
\DeclareMathOperator{\GL}{GL}
\DeclareMathOperator{\SU}{SU}
\DeclareMathOperator{\Spin}{Spin}
\DeclareMathOperator{\Sym}{Sym}
\DeclareMathOperator{\Span}{span}
\DeclareMathOperator{\CT}{CT}
\DeclareMathOperator{\conv}{conv}
\DeclareMathOperator{\Supp}{Supp}
\DeclareMathOperator{\Lie}{Lie}
\DeclareMathOperator{\core}{core}
\newcommand{\der}{\mathrm{der}}
\newcommand{\frg}{\mathfrak g}
\newcommand{\frh}{\mathfrak h}
\newcommand{\sltwo}{\mathfrak{sl}_2}
\newcommand{\so}{\mathfrak{so}}
\newcommand{\slthree}{\mathfrak{sl}_3}
\setlength{\emergencystretch}{2em}
\raggedbottom
\title{A Classification of Mathieu's Second Conjecture\\for Reductive Spherical Pairs}
\author{Christopher D. Long\thanks{\texttt{galizur@gmail.com}.}}
\date{September 22, 2026\\[2pt]\small Research draft}
\hypersetup{pdftitle={A Classification of Mathieu's Second Conjecture for Reductive Spherical Pairs},pdfauthor={Christopher D. Long},pdfsubject={Reynolds moments and nullcones on affine reductive spherical homogeneous spaces},pdfkeywords={Mathieu conjecture, spherical homogeneous space, root SL2, branching rule, nullcone, Reynolds operator}}
\begin{document}
\maketitle
\begin{abstract}
Let $G$ be a connected complex reductive algebraic group and let $H\leq G$ be a reductive closed subgroup, not necessarily connected, such that $X=G/H$ is spherical. We prove that Mathieu's Second Conjecture holds for $(G,H)$ if and only if $G_{\der}\subseteq H$, equivalently if and only if the effective semisimple action on $X$ is trivial. The same condition characterizes when the kernel of the normalized Reynolds functional is a Mathieu--Zhao subspace. When the condition fails, there are fixed $P,Q\in\C[X]$ whose pure moments all vanish and whose first marked moments are strictly positive for every positive exponent; in particular, $P$ is outside the $G$-nullcone. The proof is uniform: it combines the universal radial identity of the author's Mathieu-property paper with two classical branching rules, an explicit Weyl-even circuit on $\SL_2/N(T)$, and an elementary subdirect-product argument, rather than checking the classified families of reductive spherical subgroups one by one.
\end{abstract}
\smallskip
\noindent\textit{2020 Mathematics Subject Classification.} 14L24, 14M27, 20G20, 22E46.\\
\textit{Keywords.} Mathieu's Second Conjecture, spherical pair, nullcone, Reynolds functional, root doublet.

\section{Statement and conventions}
All algebraic groups, varieties, and representations in this paper are over $\C$. The group $G$ is connected and reductive; the reductive closed subgroup $H$ may be disconnected. A homogeneous space $X=G/H$ is \emph{spherical} if a Borel subgroup of $G$ has an open orbit. Since $H$ is reductive, $X$ is affine by Matsushima's criterion \cite{Matsushima}; the disconnected case follows as well by taking the finite quotient of $G/H^\circ$.

Write
\[
 A=\C[X],\qquad
 \hX:A\longrightarrow\C
\]
for the normalized Reynolds functional: it is the unique $G$-invariant linear functional with $\hX(1)=1$. Here $A^G=\C$ by transitivity, and rational $G$-modules are locally finite and completely reducible. In compatible compact forms $L\leq K$ of $H\leq G$, respectively,
\begin{equation}\label{eq:compact}
 \hX(f)=\int_K f(kH)\dd k=\int_{K/L}f\dd\sigma,
 \qquad \int_K1\dd k=1.
\end{equation}
We use the action $(g\cdot f)(x)=f(g^{-1}x)$. The standard compact-form, highest-weight, and root-subgroup facts used below can be found in \cite{Borel,GoodmanWallach}.

For $f\in A$, the span of its $G$-orbit is finite-dimensional. The assertion $0\in\overline{G\cdot f}$ means Zariski closure in that finite-dimensional rational module. It has the same meaning in any larger finite-dimensional submodule, since a linear submodule is closed. Set
\[
 \mathcal Z_X=\{f\in A:\hX(f^m)=0\text{ for every }m\geq1\},
 \qquad
 \mathcal N_X=\{f\in A:0\in\overline{G\cdot f}\}.
\]
Mathieu's Second Conjecture is the assertion $\mathcal Z_X\subseteq\mathcal N_X$ for reductive spherical pairs \cite[p.~264]{Mathieu}. The reverse inclusion always holds: $f\mapsto\hX(f^m)$ is an invariant homogeneous polynomial on any finite-dimensional module containing $f$.

Recall that a linear subspace $M$ of a commutative unital algebra is a Mathieu--Zhao subspace if $f^m\in M$ for every $m\geq1$ implies $Qf^m\in M$ for every fixed $Q$ and all sufficiently large $m$.

\begin{theorem}[Spherical classification]\label{thm:main}
Let $G$ be connected complex reductive, let $H\leq G$ be reductive and closed, and suppose $X=G/H$ is spherical. The following are equivalent:
\begin{enumerate}[label=\textup{(\roman*)},itemsep=2pt,topsep=4pt]
\item $(G,H)$ satisfies Mathieu's Second Conjecture, that is, $\mathcal Z_X=\mathcal N_X$;
\item $\ker\hX$ is a Mathieu--Zhao subspace of $\C[X]$;
\item $G_{\der}\subseteq H$;
\item $G_{\der}$ acts trivially on $G/H$.
\end{enumerate}
If these conditions fail, there are fixed $P,Q\in\C[X]$ such that
\begin{equation}\label{eq:negative}
 \hX(P^m)=0,\qquad \hX(QP^m)>0\qquad(m\geq1),
 \qquad 0\notin\overline{G\cdot P}.
\end{equation}
The strict inequalities refer to the scalar values of the functional, not to pointwise positivity on the complex variety.
\end{theorem}

The kernel of the action is $\core_G(H)=\bigcap_{g\in G}gHg^{-1}$. Since $G_{\der}$ is normal, $G_{\der}\subseteq H$ is equivalent to $G_{\der}\subseteq\core_G(H)$, proving the equivalence of (iii) and (iv). In this case $H$ is normal and $G/H$ is an algebraic torus, including the zero-dimensional torus.

The main new structural point is that higher-rank simple factors always expose a genuine root doublet in a nontrivial module. Once they are treated, the remaining $A_1$ factors either have a proper spherical projection of the stabilizer, or are coupled diagonally. Sections~\ref{sec:doublets} and \ref{sec:structure} make these alternatives exhaustive.

Classical classification results for reductive spherical subgroups are available; in the simple-group case the foundational list is due to Kr\"amer \cite{Kramer}. One possible approach to Theorem~\ref{thm:main} would therefore be to check the classified families individually, or alternatively to prove a spherical-localization transfer principle reducing the moment problem to rank-one pieces. The argument below takes neither route. Instead, the root-doublet lemma treats all active higher-rank simple factors uniformly, while the residual $A_1$ factors are handled directly by the normalizer circuit and the subdirect-product lemma. Thus the classification is obtained without a case-by-case appeal to subgroup lists and without a separate localization-transfer theorem.

\section{Invariant moments and descent}
\begin{lemma}[Unstable elements and fixed multipliers]\label{lem:unstable}
Let a connected reductive group $G$ act rationally by algebra automorphisms on a commutative $\C$-algebra $A$, and let $\Lambda:A\to\C$ be invariant. If $0\in\overline{G\cdot f}$, then for every fixed $Q\in A$,
\[
 \Lambda(Qf^m)=0
\]
for all sufficiently large $m$. The eventual bound can be chosen uniformly for $Q$ in any fixed finite-dimensional subspace of $A$.
\end{lemma}
\begin{proof}
The case $f=0$ is immediate. Apply the affine Hilbert--Mumford criterion \cite{Kempf} in a finite-dimensional rational submodule containing $f$. There is a cocharacter $\lambda:\Gm\to G$ with $\lim_{t\to0}\lambda(t)\cdot f=0$. Thus all $\lambda$-weights occurring in $f$ are positive integers. Let their minimum be $\delta>0$. Enlarge a fixed finite-dimensional space of multipliers to a finite-dimensional rational submodule, and let $-M$ be a lower bound for its $\lambda$-weights. Since multiplication adds weights, all weights occurring in $Qf^m$ are at least $m\delta-M$. An invariant functional annihilates every nonzero $\lambda$-weight. Hence $\Lambda(Qf^m)=0$ once $m\delta>M$.
\end{proof}

Consequently, a pair with the first two identities in \eqref{eq:negative}, even with marked nonvanishing only for infinitely many exponents, puts $P$ outside the nullcone. Also, Mathieu II implies the Mathieu--Zhao property in this setting.

\begin{lemma}[Pullback and central covers]\label{lem:pullback}
Let $\pi:G_1\twoheadrightarrow G_2$ be a surjective morphism of connected reductive groups, and let $H_i\leq G_i$ be reductive with $\pi(H_1)\subseteq H_2$. For the induced surjective morphism $\bar\pi:G_1/H_1\to G_2/H_2$,
\[
 \mathfrak h_{G_1/H_1}(f\circ\bar\pi)=\mathfrak h_{G_2/H_2}(f).
\]
Thus fixed pure/marked witnesses pull back. If $\pi:\widetilde G\to G$ is a finite central isogeny and $\widetilde H=\pi^{-1}(H)$, then $\widetilde G/\widetilde H\simeq G/H$, the normalized functionals agree, and the corresponding group orbits on the coordinate ring agree.
\end{lemma}
\begin{proof}
Composition with pullback defines a normalized $G_2$-invariant functional on $\C[G_2/H_2]$, since $\pi$ is surjective. Uniqueness of the Reynolds functional proves the identity. In the second assertion the kernel of $\pi$ lies in $\widetilde H$, giving the claimed homogeneous-space isomorphism. Surjectivity also gives equality of the group orbits. Sphericity is preserved under this isomorphism: Borels map onto Borels under a central isogeny.
\end{proof}

\section{A line and a doublet supply the moment obstruction}\label{sec:line}
This section records the homogeneous-space form of the radial mechanism. The only special moment input is \cite[Theorem~4.1]{Long}.

For $z=(z_0,z_1)\in\C^2$, put
\begin{equation}\label{eq:hopf}
 a=|z_0|^2+|z_1|^2,\qquad
 u=2z_0\overline{z_1},\qquad
 v=2z_1\overline{z_0},\qquad
 \tau=|z_0|^2-|z_1|^2,
\end{equation}
so $a^2=uv+\tau^2$, and define
\begin{equation}\label{eq:radialpair}
 \mathscr P=(a+u)\bigl(a^2v-(2a+u)\tau^2\bigr),\qquad
 \mathscr Q=u.
\end{equation}
Both are invariant under common unitary scalar multiplication of $z$. For
\begin{equation}\label{eq:cm}
 c_m=\int_0^1(1-t^2)^m\dd t=\frac{4^m(m!)^2}{(2m+1)!},
\end{equation}
the cited radial theorem states that every compactly supported finite positive $\SU(2)$-invariant measure $\mu$ on the defining module $\C^2$ satisfies
\begin{equation}\label{eq:radial}
 \int\mathscr P^m\dd\mu=0,\qquad
 \int\mathscr Q\mathscr P^m\dd\mu
 =c_m\int a^{4m+1}\dd\mu\qquad(m\geq1).
\end{equation}
The right side is positive if $\mu$ is not concentrated at zero.

\begin{proposition}[Line--doublet criterion]\label{prop:line}
Let $G$ be connected reductive and $H\leq G$ reductive. Suppose there are an irreducible rational $G$-module $V$, a nonzero $H$-stable line $\C\xi\subseteq V$, and an algebraic homomorphism $\varphi:\SL_2\to G$ such that $V$ contains a two-dimensional $\varphi(\SL_2)$-submodule $W$ isomorphic to the defining module. Then there are fixed $P,Q\in\C[G/H]$ satisfying the first two assertions in \eqref{eq:negative} for every $m\geq1$.
\end{proposition}
\begin{proof}
Choose compatible compact forms $L\leq K$. After conjugating $\varphi$ and $W$ together, but leaving $H$ and $\xi$ unchanged, assume $\varphi(\SU(2))\subseteq K$. Equip $V$ with a $K$-invariant Hermitian inner product, and let $\Pi_W$ be orthogonal projection. The map
\[
 \Phi:K\longrightarrow W\simeq\C^2,\qquad
 \Phi(k)=\Pi_W\rho(k)\xi
\]
is $\SU(2)$-equivariant. Thus $\mu=\Phi_*(\dd k)$ is compactly supported and $\SU(2)$-invariant.

It is not concentrated at zero. Otherwise, continuity and full support of Haar measure would imply $\Pi_W\rho(k)\xi=0$ for every $k$. But the span of $K\xi$ is $V$, by irreducibility and Zariski density of $K$, so this would give $V\subseteq W^\perp$. Equivalently, there is a nonempty open set on which $a(\Phi(k))>0$, and it has positive Haar measure.

The action of $H$ on $\C\xi$ is an algebraic character $\chi$. Its restriction to $L$ is unitary, so
\[
 \Phi(kl)=\chi(l)\Phi(k)\qquad(k\in K,l\in L).
\]
Phase balance makes $P=\mathscr P\circ\Phi$ and $Q=\mathscr Q\circ\Phi$ right $L$-invariant. Each coordinate of $\Phi$ is a matrix coefficient; its complex conjugate on $K$ is a coefficient of the dual representation. Hence $P,Q$ extend to regular functions on $G$. Right $L$-invariance extends to right $H$-invariance because $L$ is Zariski dense in $H$, including its components. Therefore $P,Q\in\C[G/H]$. Equations \eqref{eq:compact} and \eqref{eq:radial} give
\[
 \hX(P^m)=0,\qquad
 \hX(QP^m)=c_m\int_K a(\Phi(k))^{4m+1}\dd k>0.
\]
\end{proof}

\begin{remark}\label{rem:line}
Only an $H$-stable line is needed, not an $H$-fixed vector. This distinction is important for disconnected stabilizers and for descent through finite central covers. No assumption that $\xi$ belongs to $W$ is required.
\end{remark}

\section{Every higher-rank simple factor exposes a doublet}\label{sec:doublets}
\begin{lemma}[Universal root doublet]\label{lem:doublet}
Let $S$ be a connected complex simple algebraic group of rank at least two. Every nontrivial finite-dimensional rational $S$-module contains a defining two-dimensional submodule for some root homomorphism $\SL_2\to S$.
\end{lemma}
\begin{proof}
First consider the Lie algebra $\slthree$. An irreducible nontrivial module has highest weight $(a,b)$ with $a,b\in\Z_{\geq0}$, not both zero. Extend it to the polynomial $\GL_3$-module of highest weight $(a+b,b,0)$. The classical $\GL_3\downarrow\GL_2$ branching rule \cite[Theorem~2.1]{Molev} contains every highest weight $(\mu_1,\mu_2)$ satisfying
\[
 a+b\geq\mu_1\geq b\geq\mu_2\geq0.
\]
If $b\geq1$, take $(\mu_1,\mu_2)=(b,b-1)$; if $b=0$, take $(1,0)$. In either case $\mu_1-\mu_2=1$, so this $\GL_2$-summand restricts to the defining module of the root $\SL_2$.

Next consider type $B_2=C_2$, using $\Spin_5\supseteq\Spin_4\simeq\SL_2\times\SL_2$. A highest weight is a pair $(a,b)$ with $a\geq b\geq0$, with $a,b$ both integers or both half-integers. The branching rule \cite[Section~4.1, equation~(4.4)]{Molev} contains every $\Spin_4$ highest weight $(c,d)$ of the same parity satisfying
\begin{equation}\label{eq:b2}
 a\geq c\geq b\geq|d|.
\end{equation}
The two $\SL_2$ highest weights are $c-d$ and $c+d$. If $b>0$, then $b\geq\tfrac12$, and $(c,d)=(b,b-1)$ satisfies \eqref{eq:b2}, with $c-d=1$. If $b=0$, nontriviality gives $a\geq1$, and $(c,d)=(1,0)$ works. Thus a long-root $\SL_2$ acts on a two-dimensional summand. This includes spin representations; no integrality assumption excluding them is made.

For a general simple Lie algebra of rank at least two, choose adjacent simple roots. Their rank-two semisimple subalgebra has type $A_2$, $B_2$, or $G_2$. In type $G_2$, the long roots give a closed $A_2$ subsystem. Consequently there is a simple subalgebra of type $A_2$ or $B_2$ generated by root spaces. A nontrivial representation of a simple Lie algebra is faithful at the Lie-algebra level, so its restriction to this subalgebra is nontrivial. Complete reducibility supplies a nontrivial irreducible constituent. Apply the two cases just proved. The resulting root $\sltwo$ integrates to an algebraic root homomorphism $\SL_2\to S$, and the two-dimensional submodule is the defining group module. These assertions can equivalently be made after passage to the simply connected cover and then composed with the covering map.
\end{proof}

\begin{remark}[Weights versus submodules]\label{rem:weight}
The occurrence of the numerical weight $1$ alone does not prove the existence of a defining summand: $\Sym^3(\C^2)$ has weights $3,1,-1,-3$ and no such summand. Lemma~\ref{lem:doublet} uses actual branching constituents, not that invalid weight test.
\end{remark}

\begin{lemma}[Detecting an active factor]\label{lem:detect}
Let $G/H$ be affine with $G$ connected reductive and $H$ reductive. If a connected normal simple subgroup $S\leq G$ acts nontrivially on $G/H$, there is an irreducible rational $G$-module $V$ with $V^H\neq0$ on which $S$ acts nontrivially.
\end{lemma}
\begin{proof}
Regular functions separate points of the affine variety, so $S$ acts nontrivially on $\C[G/H]$. Choose an irreducible $G$-submodule $U$ on which it acts nontrivially, and set $V=U^*$. Evaluation at $eH$ is an $H$-fixed vector $\xi\in V$. It is nonzero: if every $f\in U$ vanished at $eH$, $G$-stability of $U$ would imply $f(gH)=0$ for every $g$, hence $U=0$. The action on the dual remains nontrivial.
\end{proof}

\begin{corollary}\label{cor:highrank}
If an active normal simple factor of $G$ has rank at least two, then $G/H$ has fixed pure/marked witnesses as in \eqref{eq:negative}. This assertion does not require sphericity.
\end{corollary}
\begin{proof}
Use Lemma~\ref{lem:detect}, Lemma~\ref{lem:doublet}, and Proposition~\ref{prop:line}. When needed, pass to a finite central cover as in Lemma~\ref{lem:pullback}.
\end{proof}

\section{The Weyl-even circuit on \texorpdfstring{$\SL_2/N(T)$}{SL2/N(T)}}\label{sec:normalizer}
Let $T\leq\SL_2$ be the diagonal torus and let $N=N_{\SL_2}(T)$. For
\[
 g=\begin{pmatrix}x&y\\z&w\end{pmatrix},\qquad xw-yz=1,
\]
define the right $T$-invariant regular functions
\[
 u=-2xy,\qquad v=2zw,\qquad \tau=xw+yz,
 \qquad uv+\tau^2=1.
\]
Right multiplication by $\left(\begin{smallmatrix}0&-1\\1&0\end{smallmatrix}\right)$ sends $(u,v,\tau)$ to $(-u,-v,-\tau)$. Hence
\[
 U=u^2,\qquad V=v^2,\qquad S=\tau^2
\]
are right $N$-invariant and satisfy $UV=(1-S)^2$. Define
\begin{equation}\label{eq:evenpair}
 \begin{split}
 R(U,S)&=S(U+2)\bigl(2-S(U^2+2U+2)\bigr),\\
 P_N&=(1+U)\bigl(V-R(U,S)\bigr),\qquad Q_N=U.
 \end{split}
\end{equation}
In particular $P_N,Q_N\in\C[\SL_2/N]$. The elementary identity
\begin{equation}\label{eq:evenidentity}
 U\bigl(V-R(U,S)\bigr)=\bigl(1-(1+U)^2S\bigr)^2
\end{equation}
is obtained by substituting $UV=(1-S)^2$.

\begin{proposition}[Normalizer coefficient identity]\label{prop:normalizer}
Let $\mathfrak h_N$ be the normalized Reynolds functional on $\C[\SL_2/N]$. For every $m\geq1$ and $B\in\C[U]$,
\begin{equation}\label{eq:evenmaster}
 \mathfrak h_N\bigl(B(U)P_N^m\bigr)
 =c_{2m}[U^m]B(U)(1+U)^{m-1}.
\end{equation}
Thus $\mathfrak h_N(P_N^m)=0$ and $\mathfrak h_N(Q_NP_N^m)=c_{2m}>0$ for every $m\geq1$.
\end{proposition}
\begin{proof}
On the compact form, write
\[
 g=\begin{pmatrix}z_0&-\overline{z_1}\\ z_1&\overline{z_0}\end{pmatrix},
 \qquad |z_0|^2+|z_1|^2=1.
\]
Then $u=2z_0\overline{z_1}$, $v=\overline u$, and $\tau=|z_0|^2-|z_1|^2$. Hopf coordinates give $\tau=t$ and $u=\sqrt{1-t^2}e^{i\theta}$, with normalized measure $\frac12\dd t\,\frac{\dd\theta}{2\pi}$ after discarding the unused common phase. Thus $U=(1-t^2)e^{2i\theta}$. For $|t|<1$, equation~\eqref{eq:evenidentity} gives
\[
 P_N=U^{-1}(1+U)\bigl(1-(1+U)^2t^2\bigr)^2.
\]
This is only a temporary localization: \eqref{eq:evenpair} is regular everywhere, so the apparent pole produces no integrability issue.

Phase averaging extracts the Laurent constant term in $U$; the result is even in $t$. Consequently
\begin{align*}
 \mathfrak h_N\bigl(B(U)P_N^m\bigr)
 &=[U^m]B(U)(1+U)^m\int_0^1\bigl(1-(1+U)^2t^2\bigr)^{2m}\dd t\\
 &=[U^m]B(U)(1+U)^{m-1}J_{2m}(1+U),
\end{align*}
where $J_{2m}(y)=\int_0^y(1-s^2)^{2m}\dd s$ is a polynomial. Moreover
\[
 J_{2m}(1+U)-J_{2m}(1)
 =\int_0^U[-y(2+y)]^{2m}\dd y
 \in U^{2m+1}\C[U].
\]
Since $B(U)(1+U)^{m-1}$ has no negative powers, the error cannot affect the coefficient of $U^m$. Replace $J_{2m}(1+U)$ by $J_{2m}(1)=c_{2m}$ to obtain \eqref{eq:evenmaster}. Taking $B=1$ and $B=U$ proves the last assertions. Integration over $\SU(2)$ agrees with normalized integration over its quotient by $N\cap\SU(2)$ for these invariant functions.
\end{proof}

The formula also gives the full marker tower
\[
 \mathfrak h_N(Q_N^sP_N^m)=c_{2m}\binom{m-1}{s-1}
 \quad(1\leq s\leq m),
\]
with zero for $s>m$. Pullback handles $\SL_2/T$ as well. Importantly, no claim is made that a witness for $T$ descends merely by squaring it; \eqref{eq:evenpair} is a different, explicitly invariant pair.

\section{The residual product of \texorpdfstring{$A_1$}{A1} factors}\label{sec:structure}
\begin{lemma}[Spherical projections]\label{lem:projection}
Suppose $G=Z\times\SL_2^r$, where $Z$ is a torus, and $H\leq G$ is reductive with $G/H$ spherical. For each factor projection $p_i:G\to\SL_2$, the image $H_i=p_i(H)$ is conjugate to one of $T$, $N(T)$, or $\SL_2$.
\end{lemma}
\begin{proof}
The image $H_i$ is a closed reductive algebraic subgroup. The morphism
\[
 G/H\longrightarrow\SL_2/H_i,\qquad gH\longmapsto p_i(g)H_i
\]
is surjective. The image of an open Borel orbit is a dense orbit of a Borel in $\SL_2$, hence open because algebraic-group orbits are locally closed. Thus $\SL_2/H_i$ is spherical and $3-\dim H_i\leq2$. In particular $H_i$ is not finite.

A proper positive-dimensional reductive subgroup of $\SL_2$ has identity component a maximal torus. It normalizes that identity component, so, after conjugation, $T\subseteq H_i\subseteq N(T)$. Since $N(T)/T$ has order two, $H_i=T$ or $N(T)$. The remaining possibility is $H_i=\SL_2$.
\end{proof}

If a projection is proper, $H_i\subseteq N(T)$ after conjugation. Therefore an actual surjective morphism
\begin{equation}\label{eq:actualquotient}
 G/H\longrightarrow\SL_2/N(T)
\end{equation}
is available, and Proposition~\ref{prop:normalizer} pulls back via Lemma~\ref{lem:pullback}.

\begin{lemma}[Subdirect stabilizers]\label{lem:subdirect}
Let $G=Z\times\SL_2^r$ and let $H\leq G$ be reductive. Suppose every projection $H\to\SL_2$ is surjective. Then either $G_{\der}\subseteq H$, or $G/H$ satisfies the line--doublet criterion of Proposition~\ref{prop:line}.
\end{lemma}
\begin{proof}
The projections of $H^\circ$ are also surjective, since their images have finite index in the connected group $\SL_2$. Write
\[
 \frh=\mathfrak z(\frh)\oplus[\frh,\frh].
\]
The center projects trivially into each $\sltwo$: its image centralizes the whole surjective image. The semisimple algebra $[\frh,\frh]$ therefore projects onto every $\sltwo$.

Decompose $[\frh,\frh]$ into simple ideals. For a fixed coordinate $i$, the image of each simple ideal is an ideal of $\sltwo$, hence zero or all of $\sltwo$. Exactly one ideal has a nonzero image: two such images would commute, which is impossible in $\sltwo$. A nonzero projection of a simple ideal is an isomorphism with $\sltwo$. Each ideal projects nontrivially somewhere, since a semisimple ideal cannot embed nontrivially in the torus algebra. Thus the coordinates are partitioned into nonempty blocks, one for each simple ideal, and the image is diagonal on each block, up to automorphisms of $\sltwo$. These automorphisms are inner. After conjugating in the individual factors, the connected derived subgroup is therefore
\begin{equation}\label{eq:blocks}
 [H^\circ,H^\circ]=\prod_{B}\Delta_B\SL_2.
\end{equation}
Equality holds at group level because connected algebraic subgroups in characteristic zero are determined by their Lie algebras. The connected center of $H^\circ$ lies in $Z$.

If every block is a singleton, \eqref{eq:blocks} contains every $\SL_2$ factor, and $G_{\der}\subseteq H$. Otherwise choose distinct $i,j$ in one block and take the irreducible $G$-module
\[
 V=\C_i^2\otimes\C_j^2,
\]
with the other factors and $Z$ acting trivially. Its $H^\circ$-fixed subspace is the one-dimensional alternating line
\[
 \C\xi,\qquad \xi=e_1\otimes e_2-e_2\otimes e_1.
\]
Since $H^\circ$ is normal in $H$, this one-dimensional space is stable under the full group $H$, including every component. The $i$th $\SL_2$ factor has a defining doublet $\C_i^2\otimes\C e_1\subseteq V$. Proposition~\ref{prop:line} applies.
\end{proof}

\begin{remark}\label{rem:sphericity}
Lemma~\ref{lem:subdirect} does not use sphericity and does not require a bound on block sizes. In the residual analysis sphericity is used only in Lemma~\ref{lem:projection}, to exclude finite projections of the stabilizer. Thus, even in this residual step, the proof does not appeal to a case-by-case classification of spherical subgroups of $\SL_2^r$ or to spherical-root localization.
\end{remark}

\section{Completion of the classification}
\begin{proof}[Proof of Theorem~\ref{thm:main}]
We first prove the positive implication. Suppose $G_{\der}\subseteq H$. The subgroup $H$ is the full inverse image of a closed subgroup of the torus $G/G_{\der}$. Hence $H$ is normal and $X=G/H$ is a connected algebraic torus. The action is the surjective translation action through $G\to X$, and
\[
 \C[X]=\C[M],\qquad \hX=\CT
\]
for the character lattice $M$ of $X$.

Let $0\neq f\in\mathcal Z_X$. The Duistermaat--van der Kallen theorem \cite{DvK} gives $0\notin\conv(\Supp f)$. Strict separation followed by rational approximation and scaling produces $\nu\in\operatorname{Hom}(M,\Z)$ with $\langle\chi,\nu\rangle>0$ for every $\chi\in\Supp f$. For the translation action on functions, the cocharacter $t\mapsto\nu(t)^{-1}$ multiplies $\chi$ by $t^{\langle\chi,\nu\rangle}$ and hence sends $f$ to zero as $t\to0$. Its translates are in $G\cdot f$, since $G\to X$ is surjective. Thus $f\in\mathcal N_X$. No lifting of the cocharacter from $X$ to $G$ is necessary. The zero polynomial and the point quotient cause no exception. This proves (iii)$\Rightarrow$(i). Lemma~\ref{lem:unstable} gives (i)$\Rightarrow$(ii).

Now assume $G_{\der}\not\subseteq H$. Take a finite central isogeny
\[
 \widetilde G=Z\times\prod_{j=1}^s S_j\longrightarrow G,
 \qquad \widetilde H=\pi^{-1}(H),
\]
where $Z$ is a torus and the $S_j$ are simply connected simple groups. By Lemma~\ref{lem:pullback}, this does not change the homogeneous space, the functional, or the relevant group orbits.

If some simple factor $S_j$ of rank at least two is not contained in $\widetilde H$, its action on $\widetilde G/\widetilde H$ is nontrivial: at the base point, trivial action would force containment. Corollary~\ref{cor:highrank} supplies $P,Q$ with the required moment identities.

Otherwise every higher-rank factor is contained in $\widetilde H$. Their product is normal in $\widetilde G$, so quotienting it out leaves an isomorphic homogeneous space with acting group
\[
 G'=Z\times\SL_2^r
\]
and a reductive stabilizer $H'$. Sphericity is unchanged, because the omitted factors act trivially. The derived subgroup of $G'$ is still not contained in $H'$; if it were, all simple factors of $\widetilde G$ would be contained in $\widetilde H$.

Apply Lemma~\ref{lem:projection}. If a stabilizer projection is proper, use the genuine quotient \eqref{eq:actualquotient} and pull back the normalizer witnesses of Proposition~\ref{prop:normalizer}. If every projection is surjective, Lemma~\ref{lem:subdirect} gives a line and a doublet: its alternative $G'_{\der}\subseteq H'$ is excluded. Proposition~\ref{prop:line} again gives the witnesses. Pull them back through the homogeneous-space isomorphisms to obtain, on the original $X$,
\[
 \hX(P^m)=0,\qquad \hX(QP^m)>0\qquad(m\geq1).
\]
The Mathieu--Zhao property fails. Lemma~\ref{lem:unstable} also gives $0\notin\overline{G\cdot P}$, so Mathieu II fails. This proves both the remaining equivalences and the stronger assertion.
\end{proof}

\begin{corollary}[Group-type specialization]
For every connected complex reductive group $G$, the pair $(G\times G,\Delta G)$ satisfies Mathieu's Second Conjecture if and only if $G$ is an algebraic torus.
\end{corollary}
\begin{proof}
The pair is spherical by algebraic Peter--Weyl theory. Its derived subgroup is $G_{\der}\times G_{\der}$, which is contained in $\Delta G$ exactly when $G_{\der}=1$. Apply Theorem~\ref{thm:main}.
\end{proof}

\begin{remark}[Scope and proof dependencies]
The theorem treats every reductive spherical pair with connected complex acting group, including disconnected stabilizers, nontrivial centers, and finite central forms. It does not claim a classification for nonspherical homogeneous spaces, disconnected acting groups, or positive characteristic. The special external moment input is \cite[Theorem~4.1]{Long}, not that paper's classification theorem. The additional representation-theoretic inputs are the two branching rules used in Lemma~\ref{lem:doublet}; the residual subgroup argument and the normalizer circuit are proved here.
\end{remark}

\paragraph{Draft preparation.}
ChatGPT (OpenAI) assisted with the proof development, drafting, source checking, and symbolic checks. This research draft has not been independently peer reviewed or formally verified.

\begin{thebibliography}{10}
\bibitem{Borel}
A.~Borel, \emph{Linear Algebraic Groups}, second ed., Graduate Texts in Mathematics 126, Springer, 1991.
\href{https://doi.org/10.1007/978-1-4612-0941-6}{doi:10.1007/978-1-4612-0941-6}.
\bibitem{DvK}
J.~J.~Duistermaat and W.~van der Kallen, \emph{Constant terms in powers of a Laurent polynomial}, Indag. Math. (N.S.) \textbf{9} (1998), no.~2, 221--231.
\href{https://doi.org/10.1016/S0019-3577(98)80020-7}{doi:10.1016/S0019-3577(98)80020-7}.
\bibitem{GoodmanWallach}
R.~Goodman and N.~R.~Wallach, \emph{Symmetry, Representations, and Invariants}, Graduate Texts in Mathematics 255, Springer, 2009.
\href{https://doi.org/10.1007/978-0-387-79852-3}{doi:10.1007/978-0-387-79852-3}.
\bibitem{Kempf}
G.~R.~Kempf, \emph{Instability in invariant theory}, Ann. of Math. (2) \textbf{108} (1978), no.~2, 299--316.
\href{https://doi.org/10.2307/1971168}{doi:10.2307/1971168}.
Earlier version, transcribed by I.~Morrison:
\href{https://arxiv.org/abs/1807.02890}{arXiv:1807.02890}.
\bibitem{Kramer}
M.~Kr\"amer, \emph{Sph\"arische Untergruppen in kompakten zusammenh\"angenden Liegruppen}, Compositio Math. \textbf{38} (1979), no.~2, 129--153.
\href{https://www.numdam.org/item/CM_1979__38_2_129_0/}{Numdam}.
\bibitem{Long}
C.~D.~Long, \emph{The Mathieu Property for Compact Connected Lie Groups}, preprint, 2026.
\href{https://arxiv.org/abs/2609.16178v1}{arXiv:2609.16178v1}.
\bibitem{Mathieu}
O.~Mathieu, \emph{Some conjectures about invariant theory and their applications}, in \emph{Alg\`ebre non commutative, groupes quantiques et invariants} (Reims, 1995), S\'eminaires et Congr\`es 2, Soci\'et\'e Math\'ematique de France, 1997, 263--279.
\href{https://smf.emath.fr/publications/some-conjectures-about-invariant-theory-and-their-applications}{SMF publication record}.
\bibitem{Matsushima}
Y.~Matsushima, \emph{Espaces homog\`enes de Stein des groupes de Lie complexes}, Nagoya Math. J. \textbf{16} (1960), 205--218.
\href{https://doi.org/10.1017/S0027763000007662}{doi:10.1017/S0027763000007662}.
\bibitem{Molev}
A.~I.~Molev, \emph{Gelfand--Tsetlin bases for classical Lie algebras}, in \emph{Handbook of Algebra}, vol.~4, Elsevier, 2006, 109--170.
\href{https://doi.org/10.1016/S1570-7954(06)80006-9}{doi:10.1016/S1570-7954(06)80006-9}.
\href{https://arxiv.org/abs/math/0211289v2}{arXiv:math/0211289v2}; theorem and equation numbering cited from this version.
\end{thebibliography}
\end{document}